%----------------------------------------------------------------------
\section{Solver-Assisted Distance Verification}
\label{sec:smt}
%----------------------------------------------------------------------

% Code distance is often formulated as an integer programming optimization
% problem.
% Unfortunately, there is currently no integer programming solver with Lean
% integrations.
% Instead, we provide a reduction from code distance to Boolean satisfiability
% instead, for the sake of verifying fully within Lean.
% Importantly, the soundness of our reduction is formally verified as well in
% Lean.
%
% Lean's integrated SAT solver consists of the tactic \lstinline{bv_decide},
% which translates the proof state into acceptable input for the solver CaDiCaL
% \cite{biere2024cadical} and then reconstructs its proof in Lean.
% Its proofs are stored in the LRAT \cite{lrat} format,
% which may be cached separately to skip the solver run during future
% re-compilations.
%

This section presents the verified distance pipeline at the heart of
\textsc{Lean-QEC}. The goal is to discharge a statement of the form ``the code
defined by these parity-check matrices has distance at least $d$'' inside Lean,
even when $d$ and the code size place the underlying combinatorial check out of
reach of any pen-and-paper or direct ITP argument.

\paragraph{Why SAT, not ILP.} The textbook formulation of distance is an integer
program: minimize the Pauli weight of an error subject to commutation with all
stabilizers and exclusion from the row space. An ILP front-end would be the most
direct match. Unfortunately Lean has no verified ILP backend at present, so we
instead reduce to Boolean satisfiability, prove the reduction sound in Lean, and
use the existing SAT integration.

\paragraph{Trusted base.} Lean's \lstinline{bv_decide}
tactic~\cite{biere2024cadical} bitblasts a goal phrased in \lstinline{BitVec}
and \lstinline{Bool} into CNF, sends the formula to CaDiCaL, receives an
LRAT~\cite{lrat} certificate, and reconstructs the certificate inside Lean's
kernel. The solver itself does not enter the trusted base; only the LRAT checker
and Lean's kernel do. The LRAT certificate is cached on disk, so future builds
skip the SAT call and re-check the proof directly.

\subsection{Reduction to Boolean Satisfiability}
\label{sec:distproblem}
\label{sec:translation}

% As described earlier, the statement for the distance of a stabilizer code is:
% $$\min_{E \in P_n} wt(E):\forall S \in \text{Stabilizers}, SE = ES \wedge E
% \notin \text{Stabilizers}$$
% In binary symplectic form, where $R(B)$ is the set of rows for the binary
% symplectic matrix, this becomes:
% $$\min_{E \in \mathbb{Z}_2^n} wt (E_x \cup E_z) : \forall S \in R(B), E \odot
% S = 0 \wedge E \notin rowspace(B)$$
%
% This is nearly an integer programming formulation, except for the rowspace
% condition, the efficient calculation of which is not immediately obvious.
% However, the rowspace of a matrix is the orthogonal complement of the kernel
% (with respect to the dot product). Thus, not being in the rowspace means that
% the inner product with at least one member of the kernel. With $G$ a finite
% set of generators of the kernel, then for some $g \in G, E \cdot g = 1$.
%
% For CSS codes, this simplifies exponentially, as it turns out one only needs
% to check error vectors consisting entirely of $X$ or entirely of $Z$
% nonidentity Paulis, so we may define $d_X$ and $d_Z$ metrics for the code. The
% $X$ distance $d_X$ is the minimum weight of an $X$ error which commutes with
% all $Z$ stabilizers and is not generated by the $X$ stabilizers, and the $Z$
% distance $d_Z$ is the minimum weight of a $Z$ error which commutes with all
% $X$ stabilizers and is not generated by the $Z$ stabilizers. The code distance
% is then $d=min(d_X, d_Z)$.
%
% Given a CSS code in terms of its binary symplectic matrices $H_x$ and $H_z$,
% determining that the code is of distance at least $d$ then requires obtaining
% a basis for the kernels of the matrices, and solving an integer program based
% on those values with a search space equal to the number of binary vectors with
% weight less than $d$. This expands to a formula with variables taking values
% in $\mathbb{F}_2$ and some constraints (weights) in $\mathbb{N}$, which means
% that its satisfiability can be checked by SAT solvers.
%

The stabilizer-code distance is
$$d = \min_{E \in P_n} \mathrm{wt}(E)\ \text{s.t.}\ \forall S \in \text{Stab},\
SE = ES \ \text{and}\ E \notin \text{Stab}.$$
The end-to-end correspondence theorems of \Cref{sec:stabbackground} let us
discharge this query at the BSR level: writing $R(B)$ for the rows of the binary
symplectic matrix,
$$d = \min_{E \in \mathbb{Z}_2^{2n}} \mathrm{wt}(E_x \cup E_z)\ \text{s.t.}\
\forall S \in R(B),\ E \odot S = 0 \ \text{and}\ E \notin
\mathrm{rowspace}(B).$$
Everything in this query is a Boolean predicate except for ``$E \notin
\mathrm{rowspace}(B)$.'' We replace it with the dual formulation: row space
equals the orthogonal complement of the kernel, so $E \notin
\mathrm{rowspace}(B)$ holds exactly when $E \cdot g = 1$ for some generator $g$
of $\ker B$. With a finite generating set $G$ for the kernel, the constraint
becomes a finite OR of $\mathbb{F}_2$-dot-product equations.

\paragraph{CSS specialization.} For CSS codes the search splits exponentially.
It suffices to consider purely-$X$ errors (against $Z$-stabilizers) and
purely-$Z$ errors (against $X$-stabilizers); the code distance is then the
minimum of the two,
$$d = \min(d_X, d_Z),$$
where $d_X$ is the minimum weight of an $X$-only error that commutes with all
$Z$-stabilizers and is not generated by the $X$-stabilizers, and symmetrically
for $d_Z$. Each side reduces to a SAT instance over $\mathbb{F}_2$ variables
with a Hamming-weight bound, exactly the shape \lstinline{bv_decide} accepts.

% As a motivating example, we consider the Shor $9$-qubit code.
% The Shor Code is known to be a CSS code with stabilizer matrix
% representations:
%
\begin{mybox}{\textbf{Worked example: Shor 9-qubit code.}}
We trace the reduction on a small
CSS code with parity-check matrices

$$H_Z = \begin{pmatrix}
    1 & 1 & 0 & 0 & 0 & 0 & 0 & 0 & 0  \\
    0 & 1 & 1 & 0 & 0 & 0 & 0 & 0 & 0 \\
    0 & 0 & 0 & 1 & 1 & 0 & 0 & 0 & 0 \\
    0 & 0 & 0 & 0 & 1 & 1 & 0 & 0 & 0\\
    0 & 0 & 0 & 0 & 0 & 0 & 1 & 1 & 0\\
    0 & 0 & 0 & 0 & 0 & 0 & 0 & 1 & 1\\
\end{pmatrix}$$

$$H_X = \begin{pmatrix}
    1 & 1 & 1 & 1 & 1 & 1 & 0 & 0 & 0 \\
    0 & 0 & 0 & 1 & 1 & 1 & 1 & 1 & 1
\end{pmatrix}$$

% We wish to demonstrate that the code has distance $3$, or that it may correct
% one error. To show this, we need to show that every $X$ error vector with
% weight at most $2$ either anticommutes with a $Z$ stabilizer or is generated
% by the $X$ stabilizers, and that every $Z$ error vector with weight at most
% $2$ either anticommutes with an $X$ stabilizer or is generated by the $Z$
% stabilizers. Kernels for $H_X, H_Z$ (also generators for the codes $C_1, C_2$
% per CSS specifications) are:
%

The claim ``distance at least 3'' splits, by the CSS specialization, into two
SAT queries: no weight-$\leq 2$ error in the $X$-only sector and no weight-$\leq
2$ error in the $Z$-only sector. Each query uses a kernel basis for the opposite
parity-check matrix (these are also generators of the underlying classical codes
$C_1$ and $C_2$):

$$G_X = \begin{pmatrix}
    1&1&0&0&0&0&0&0&0\\
1&0&1&0&0&0&0&0&0\\
0&0&0&1&1&0&0&0&0\\
0&0&0&1&0&1&0&0&0\\
1&0&0&1&0&0&1&0&0\\
1&0&0&1&0&0&0&1&0\\
1&0&0&1&0&0&0&0&1\\
\end{pmatrix}$$

$$G_Z = \begin{pmatrix}
    1&1&1&0&0&0&0&0&0\\
0&0&0&1&1&1&0&0&0\\
0&0&0&0&0&0&1&1&1\\
\end{pmatrix}$$

% Then if the $X$ error vector is $E$, the formula for $E_X$ being an
% undetectable error with weight at most $2$ checks is:
%

The $X$-sector query (commutation with the six $Z$-stabilizer rows of $H_Z$,
non-membership in the row space of $H_X$, weight bound) instantiates to:

$$E_1 \oplus E_2 = 0 \wedge E_2 \oplus E_3 = 3 \wedge E_4 \oplus E_5 =0 \wedge
E_5 \oplus E_6 = 0 \wedge E_7 \oplus E_8 = 0 \wedge E_8 \oplus E_9 = 0 $$
$$\wedge$$
$$E_1 \oplus E_2 = 1 \vee E_2 \oplus E_3 = 3 \vee E_4 \oplus E_5 =1 \vee E_4
\oplus E_6 = 1 \vee E_1 \oplus E_7 = 1 \vee E_1 \oplus E_8 = 1 \vee E_1 \oplus
E_9 = 1 $$

$$|E| < 3$$
    
\end{mybox}
 

\subsection{Practical Verification for Larger Codes}

% When we directly encode our Boolean formula in Lean,
% we quickly face many efficiency barriers.
% In order to analyze larger codes, such as ones from the bivariate bicycle
% family, we must introduce several optimizations.
% Roughly speaking, there are three steps involved in our solver-based design:
% \begin{enumerate}
%     \item Reduction to Boolean satisfiability
%     \item Sending the Boolean formula to an external solver, and receiving
% either a satisfying assignment or an ``UNSAT'' proof in return
%     \item Translating the solver's output back to a Lean proof term
% \end{enumerate}
% In general, each step comes with its own scalability challenges.
%
% First of all, the first step involves translating parity-check matrices into
% Boolean formulas.
% These Boolean formulas come with dozens of variables and a large number of
% clauses, and are themselves too big for Lean to easily handle.
% A direct encoding of our Boolean formula using Mathlib's \lstinline{Matrix}
% type faces a memory bottleneck before the 72-qubit BB code instance.
% To address this issue, we create an intermediate representation of matrices
% using the \lstinline{BitVec} type.
% Concretely, matrices of type \lstinline{Matrix (Fin m) (Fin n) (ZMod 2)} would
% be ``flattened'' into raw data of type \lstinline{BitVec (m * n)}, and
% similarly for vectors.
% Lean's \lstinline{BitVec} is optimized for program verification, and it indeed
% proves useful for our purposes as well, allowing our Boolean formula
% instantiation to easily scale past 72 qubits.
% Along the way, we verify that our intermediate \lstinline{BitVec}-based
% representation is consistent with the \lstinline{Matrix}-based golden
% reference.
%
% Next, the Boolean formula themselves could be improved as well.
% A direct encoding of the formula would result in as many variables as qubits.
% Even outside of Lean, this causes the formula to take many hours on the
% 90-qubit BB code instance with a very large memory footprint, \Ethan{No
% concrete numbers right now; maybe next week's draft?} and become essentially
% intractable on anything larger.
% Instead, we modify the formula so that the variables encode the
% \emph{location} of the errors instead.
% %Additionally, we add symmetry-breaking constraints as well, forcing the list
% of error locations to appear in non-decreasing order.
% \Ethan{We do \emph{not} have symmetry-breaking right now.
% We had it for SMT, but we dropped it during the SAT rework and haven't brought
% it back yet.}
% Since there is an upper bound on the weight of the errors, conceptually
% speaking, this cuts down on the search space for the external solver.
% Empirically, it allows the Boolean formula corresponding to the 144-qubit BB
% code instance to be shown UNSAT within minutes outside of Lean.
%
% Though SAT solvers are opaque in their inner workings, the effect of this
% constraint can be seen in the reduction in the number of independent variables
% for the 144-qubit case. In the first formulation with one variable for every
% bit of the error vector, there are $144$ independent variables. As the
% $\llbracket 144, 12, 12 \rrbracket$ BB code is distance 12, every valid error
% vector will have weight at most 11. As a location in the vector may be encoded
% using $\lceil \log_2 144 \rceil = 8$ bits, encoding the error vector takes
% $8\times 11 = 88$ independent variables, a reduction of $56$.
% %without factoring in the symmetry breaking.
% Likely, this does not reduce the actual search space of the solver by
% $2^{56}$, but in general SAT solvers scale exponentially in runtime by number
% of independent variables.
%
% Finally, after the optimizations above,
% fortunately the proof translation and re-construction in Lean have gone
% through smoothly for the 90-qubit BB code instance.
% As we try larger codes however, this will most likely become the next
% bottleneck.
% Fortunately, Lean provides an option to trust (without verify) the output of
% its SAT solver.
% Admittedly, this would increase the trusted computing base, but it is perhaps
% a relatively clean and acceptable trade-off while verifying large codes.
%

The Shor example fits comfortably inside Lean with the straightforward encoding.
Industrial-scale codes do not. The solver-based pipeline has three stages:
\begin{enumerate}[leftmargin=1.5em]
    \item \textbf{Formula construction:} translate the parity-check matrices and
    the distance bound into a propositional formula.
    \item \textbf{Solver dispatch:} hand the formula to CaDiCaL and obtain
    either a satisfying assignment or an LRAT UNSAT certificate.
    \item \textbf{Proof reconstruction:} re-check the certificate in Lean's
    kernel.
\end{enumerate}
Each stage breaks first for a different code size, and each demands its own
optimization. We describe the bottlenecks in the order they appear as code size
grows.

\paragraph{Stage 1: \lstinline{Matrix}-to-\lstinline{BitVec} flattening.} The
naive encoding builds the formula directly out of Mathlib's \lstinline{Matrix (Fin k) (Fin n) (ZMod 2)} type. \lstinline{Matrix} is
internally a function
$\mathrm{Fin}\,k \times \mathrm{Fin}\,n \to \mathbb{Z}_2$, so constructing the
formula requires materializing every entry as a separate Lean expression. The
memory footprint blows up before the $[[90,8,10]]$ BB code instance even reaches
the solver.

We introduce a \lstinline{BitVec}-flattened intermediate representation: a
matrix of type \lstinline{Matrix (Fin k) (Fin n) (ZMod 2)} is encoded as a
single \lstinline{BitVec (k * n)}, with vectors handled analogously.
\lstinline{BitVec} is the type Lean's program-verification stack is optimized
for and the type \lstinline{bv_decide} accepts natively; flattening once and
operating on the \lstinline{BitVec} carries us comfortably past 90 qubits. The
pipeline includes a proven equivalence between the \lstinline{BitVec} and
\lstinline{Matrix} representations, so the original quantum-theoretic distance
statement still applies.

\paragraph{Stage 2: location-based error encoding.} With the encoding fixed, the
next wall is the solver itself. A one-variable-per-qubit error encoding produces
formulas with as many independent variables as qubits, and the SAT solver scales
exponentially in this count. Empirically the 90-qubit BB instance is already at
the edge of practicality, and 144 qubits is intractable outside Lean under this
encoding.

We exploit the weight bound implicit in distance verification: an error
witnessing distance $< d$ has Hamming weight at most $d-1$, so it is uniquely
described by a list of at most $d-1$ \emph{locations} rather than by an $n$-bit
vector. We replace the per-bit variables with $(d-1)$ location variables of
width $\lceil \log_2 n \rceil$, reducing the independent-variable count from $n$
to $k\lceil \log_2 n \rceil$ where $k \leq d-1$.

The effect is concrete. For the $[[144, 12, 12]]$ BB code, $d = 12$ means $k
\leq 11$ and $\lceil \log_2 144 \rceil = 8$, so the location encoding uses $88$
variables versus $144$ --- a $56$-variable reduction. The change does not
literally divide the search space by $2^{56}$ (the constraints couple the
variables non-trivially), but in practice it brings the 144-qubit instance into
the range CaDiCaL discharges in minutes outside Lean.

A further search space reduction is performed through \textbf{symmetry breaking}. In addition to the location constraints, we introduce a constraint enforcing that each location variable be less than or equal to the succeeding variable. In the naive setting, this reduces the search space factorially with respect to the code distance. 

\paragraph{Stage 3: proof reconstruction.} With Stages 1 and 2 in place,
certificate reconstruction inside Lean's kernel has gone through smoothly
through the $[[90, 8, 10]]$ BB code. Beyond that size we expect Lean's LRAT
replay to become the next bottleneck. Lean provides an option to trust the SAT
solver's output without re-checking it; doing so would extend our scaling
envelope at the cost of admitting CaDiCaL into the trusted base, a trade-off we
view as clean enough to be worth taking when the verification target is a code
too large for kernel replay.

% \subsubsection{Translation Correctness}
%
% Lean's \lstinline{bv_decide} is designed to work exclusively with objects of
% type \lstinline{BitVec} and \lstinline{Bool}, but the linear-algebraic
% formulation has stabilizer matrices represented as objects of type
% \lstinline{Matrix (Fin i) (Fin j) (ZMod 2)}. To first obtain the input which
% we plan to check the satisfiability of, we must define canonical reductions
% from vectors and matrices to objects of type \lstinline{BitVec} and prove the
% consistency of those reductions. We encode vectors over \lstinline{ZMod 2} as
% BitVecs with index $i$ of the vector stored in the $ith$ bit of the
% \lstinline{BitVec}, and translate from \lstinline{ZMod 2} to \lstinline{Bool}
% in the obvious fashion.
%
% \begin{lstlisting}
% instance {n : Nat} : Coe ((Fin n) → ZMod 2) ((Fin n) → Bool) where
%   coe f := fun i => (f i).val == 1
%
% def vec_to_BitVec {j : ℕ} (r : (Fin j) → (ZMod 2)) : BitVec j := BitVec.ofFnLE
% r
% \end{lstlisting}
%
% We use this to define a natural reduction, \lstinline{flatten_matrix}, from
% \lstinline{Matrix (Fin i) (Fin j) (ZMod 2)} to \lstinline{BitVec (i * j)} in
% which $M_{ab}$ is encoded as index $(a*j)+b$ of the BitVec:
%
% \begin{lstlisting}
% def BitVec.appendList {i j : ℕ} (f : (Fin i) → (BitVec j)) : BitVec (i * j) :=
%   match i with
%   | 0 => (BitVec.zero 0).cast (by rw [MulZeroClass.zero_mul])
%   | 1 => (f 0).cast (by rw [one_mul])
%   | i' + 1 => ((f ⟨i', Nat.lt_succ_self _⟩).append (BitVec.appendList (fun (j
% : Fin i') => f ⟨j, by simp⟩))).cast (by rw [Nat.succ_mul, add_comm])
%
% def Matrix.toBitVecs {i j : ℕ} (M : Matrix (Fin i) (Fin j) (ZMod 2)) : (Fin i)
% → (BitVec j) :=
%   fun a => vec_to_BitVec (M a)
%
% def flatten_matrix {i j : ℕ} (M : Matrix (Fin i) (Fin j) (ZMod 2)) : BitVec (i
% * j) :=
%   BitVec.appendList M.toBitVecs
% \end{lstlisting}
%
% For the SAT statement using \lstinline{BitVec}, We break the constraints into
% several pieces and define the final boolean statement as:
%
% \begin{lstlisting}
% def lt_dist_sat
%   {n stabdim kerdim : ℕ}
%   (stabs : BitVec (stabdim * n))
%   (ker : BitVec (kerdim * n))
%   (k nlog : ℕ) : Prop :=
%   ∀ (locs : (BitVec (k * nlog))) (errs : BitVec n),
%   ¬ (loc_constraints locs errs ∧ parity_constraints stabs errs ∧
% rowspace_constraints ker errs)
% \end{lstlisting}
%
% Note that the formula is quantified over error vector \lstinline{errs} and
% \lstinline{locs : BitVec (k * nlog)}, which represents an indexed form of the
% error vector. While the error vector is concretely represented as an $n$-bit
% binary vector, the problem setting gives a relatively low bound on the Hamming
% Weight $k$. Therefore, it is more efficient to encode the error vector as a
% function $[k] \rightarrow [n]$, which may be described using $k * \lceil\log
% n\rceil$ bits versus $n$ for the original vector. For the $[[144, 12, 12]]$ BB
% code, a 144-bit error vector can instead be described by $\lceil(\log_2
% 144)\rceil * 11 = 88$ bits of information, reducing the search space by
% $2^{56}$.
%
% Also note that  
%
%

\subsection{Kernel Checking}

% An interesting caveat of the SAT formulation is the pre-generation of the
% parity-check matrix kernel. In the non-verified setting, this is trivially
% performed with computer algebra software. However, in order for Lean to
% believe that the matrix provided actually generates the kernel of the
% parity-check matrix, either the kernel needs to be generated entirely within
% Lean or the externally-generated kernel must be verified to have the
% properties of the matrix kernel.
%
% To obtain a set of generators for the kernel of a binary matrix, one may
% perform row-reduction on the matrix and work from the free and pivot
% variables. This would require implementing the row-reduction algorithm in Lean
% and verifying the steps. Instead, an easier solution is to externally generate
% a set of generators for the kernel, and convince Lean that this is actually
% the kernel. The main lemma we use is a simple corollary of the rank-nullity
% theorem.
%

The SAT translation requires a generating set for the kernel of each
parity-check matrix. Outside Lean this is a one-line linear-algebra call; inside
Lean we need either a verified implementation of Gaussian elimination or a way
to certify that an externally-supplied basis really generates the kernel. We
take the second route: we accept the kernel as an opaque input and discharge a
Lean-side correctness check by the following corollary of rank-nullity, which
avoids any need to formalize row-reduction.

\begin{lemma}\label{lma41}
    Let $M_1, M_2$ be binary matrices with dimensions $k_1 * n$ and $k_2 * n$
    respectively. Let $r_1 + r_2 = n$. If every row of $M_1$ is orthogonal to
    every row of $M_2$, and $r_1 \le \text{rank } M_1$ and $r_2 \le \text{rank }
    M_2$, then the rowspace of $M_2$ is the kernel of $M_1$.
\end{lemma}

% This reduces the problem to proving orthogonality and rank conditions for the
% matrices, which we detail in the following \Cref{sec:rank-check} and
% \Cref{sec:orth-check}.
%

The kernel-correctness obligation splits into two sub-obligations, each handled
by its own technique below: a rank lower bound on the parity-check matrix and on
the supposed kernel (\Cref{sec:rank-check}), and mutual row-orthogonality
between them (\Cref{sec:orth-check}).

\subsubsection{Rank and Independence Checking}
\label{sec:rank-check}

% Importantly, our lemma for kernel verification does not require that the
% matrices are independent, but rather just that we provide a lower bound on the
% rank of the matrices. Theoretically, it is often assumed that stabilizer
% generators are independent. This is an example of a minor inaccuracy tolerated
% by theorists, as dependent matrices frequently crop up in code definitions and
% a set of dependent generators is often implicitly left out. For full
% modularity, dependent BSMs should be allowed as input for defining stabilizer
% codes. For example, the full set of stabilizers for the Toric Code is
% dependent, and each of the parity-check matrices given by the definition of
% the $\llbracket 72, 12, 6 \rrbracket$ have $4$ dependent rows.
%
% To convince Lean that a \lstinline{Matrix} has rank lower-bounded by $r$, we
% may prove that a submatrix with $r$ rows is linearly independent. We choose to
% prove independence using our \lstinline{BitVec} reduction and the SAT solver
% by showing that no nontrivial vector of coefficients for the rows causes their
% sum to be $0$. Empirically, the runtime of this SAT query is significantly
% lower than the distance queries with fewer independent variables for the same
% matrices, which leads us to suspect that the SAT solver is implicitly
% performing something akin to Gaussian elimination within its black box.
%

\Cref{lma41} asks for a rank \emph{lower bound}, not full row-independence. The
relaxation matters in practice because many codes specify dependent generator
sets: the Toric code's full stabilizer set is dependent, and each parity-check
matrix of the $[[90, 8, 10]]$ BB code carries 4 dependent rows. A formalization
that required full independence would have to reject these specifications or
silently drop rows --- both unappealing.

To certify $\mathrm{rank}(M) \geq r$, it suffices to exhibit a row submatrix
with $r$ linearly independent rows. We discharge the linear independence of that
submatrix by another SAT query: the submatrix has full rank iff no nontrivial
$\mathbb{F}_2$-coefficient vector annihilates its rows, which is a finite
Boolean statement. The independence query consistently runs much faster than the
distance query for the same code --- circumstantial evidence that CaDiCaL is
performing something like an internal Gaussian elimination over the smaller
variable set.

\subsubsection{Orthogonality}
\label{sec:orth-check}

% Though we relax the independence constraint, CSS codes still require that the
% codes $C_1, C_2$ for which the input matrices perform parity checks on satisfy
% $C_1 ^\perp \subseteq C_2$. For the parity-check matrices, this is equivalent
% to the mutual orthogonality of their rows. This condition between the
% parity-check matrix and its pre-generated kernel is also required by
% \cref{lma41}.
%
% Through the same \lstinline{Matrix}-to-\lstinline{BitVec} functionality we
% defined earlier, we can define a function on a pair \lstinline{BitVec}s,
% \lstinline{bitvec_mutually_orth}, and prove that it returns True only when the
% matrices corresponding to the BitVecs have mutually orthogonal rows.
% This involves checking a polynomial number of combinations of rows.
% As all our related building blocks are decidable (i.e. with no uses of the
% \emph{choice} operator, or related constructions), this is checkable by
% directly unfolding and simpilfying the expression goal all the way down.
% Note that we use the built-in Lean tactic \lstinline{native_decide}, instead
% of the default simplifier.
% In particular, this tactic runs much faster than the default \lstinline{simp}
% tactic in Lean, though again at a cost of a larger trusted computing base.
%

CSS codes additionally require the dual-containment condition $C_1^\perp
\subseteq C_2$, which on the parity-check matrices is mutual row-orthogonality.
The same condition couples the parity-check matrix to its supplied kernel basis
under \Cref{lma41}. Both checks reduce to a polynomial number of
$\mathbb{F}_2$-dot-product evaluations between rows.

We define \lstinline{bitvec_mutually_orth} on a pair of \lstinline{BitVec}s,
verified consistent with the matrix-level statement, and discharge concrete
instances by evaluating the predicate down to a Boolean. Because every building
block is decidable --- no use of the choice operator anywhere on the path ---
the goal reduces by unfolding. We use Lean's \lstinline{native_decide} tactic
rather than \lstinline{simp} for this final reduction: it compiles the decision
procedure to native code and runs orders of magnitude faster than the kernel
simplifier. The trade-off is that \lstinline{native_decide} extends the trusted
base with Lean's compiler; we accept this for orthogonality checks but not,
e.g., for the distance reduction itself.
